\documentclass[../fractals_everywhere.tex]{subfiles}
\begin{document}
\section{Class 02 - April 17th, 2026}
\subsection{Motivations}
\begin{itemize}
	\item The Hausdorff Distance.
\end{itemize}
\subsection{The Hausdorff Distance}
\begin{exr}
	Let \((X, d)\) be a complete metric space, A and B be points in \(\mathcal{H}(X)\) such that \(A\neq B\). Show that either \(d(A,B)\neq 0\) or \(d(B, A) \neq 0\), and that if \(A\subseteq B\), then \(d(A, B) = 0.\)

	For the first part, if A is different from B, there is at least one point p in A that is not in B; we further know that if we pick the distance \(d(p, B),\) then \(d(p, B)\neq 0\), hence by definition of the distance from sets, it follows that \(\max\limits_{x\in A}(d(x, B))\geq d(p, B) > 0\). In conclusion, \(d(A, B)\neq 0\) and the same works for \(d(B, A)\).

	Now if \(A\subseteq B\), this means that any \(p\in A\) is also in B, and if we take the minimum distance of a point x in A and a subset B, then
	\[
		d(x, B) = \min\limits_{y\in B}(d(x, y)) = d(x, x) = 0,
	\]
	from which follows that \(d(A, B) = 0.\)
\end{exr}
\begin{exr}
	Let \((X, d)\) be a copmlete metric space. Show that if A, B and C are all in \(\mathcal{H}(X)\), then \(B\subseteq C\) implies \(d(A, C)\leq d(A, B)\).

	Because B is a subset of C, we have \(d(B, C) = 0\), and using the triangular inequality (tbp), it follows that
	\[
		d(A, C) \leq d(A, B) + \underbrace{d(B, C)}_{0},
	\]
	therefore \(d(A, C)\leq d(A, B).\)
\end{exr}
\begin{exr}
	Let \((X, d)\) be a copmlete metric space; show that if A, B and C are in \(\mathcal{H}(X)\), then
	\[
		d(A\cup B, C) = d(A, C) \vee d(B, C),
	\]
	where \(\cdot \vee \cdot \) denotes the maximum between two real numbers.

	To see that it is true, notice that
	\[
		d(A\cup B, C) = \max\limits_{x\in A\cup B}(d(x, C)) = \max\limits_{x\in A}(d(x, C)) \vee \max\limits_{y\in B}(d(y, C)).
	\]
	Therefore,
	\[
		d(A\cup B, C) = d(A, C)\vee d(B, C).
	\]
\end{exr}
\begin{exr}
	Let \((X, d)\) be a complete metric space, and A, B and C be sets in \(\mathcal{H}(X)\). Show that
	\[
		d(A, B) \leq d(A, C) + d(C, B).
	\]

	Pick x in A, y in C such that \(d(x, y) = d(x, C)\), and z in B such that \(d(y, z) = d(y, B)\).
	Then, we have
	\[
		d(x, z)\leq d(x, y) + d(y, z).
	\]
	By the definition of the points chosen, it also follow that
	\[
		d(x, z) \leq d(x, C) + d(y, B),
	\]
	so when looking at \(\max\limits_{x\in A}(d(x, B))\), and by the way the distance \(d(x, B)\) is defined, which is the minimum,
	\[
		d(x, B)\leq d(x, z)\leq d(A, C) + d(C, B).
	\]
	Therefore, since this is true for all x's in A,
	\[
		d(A, B)\leq d(A, C) + d(C, B).
	\]

\end{exr}

As shown by the previous exercises, \(d(A, B)\) for any two compact sets A and B is not even a metric, but we want to turn \(\mathcal{H}(X)\) into a space we can work with, so it must at least have a metric.
For that reason, we define the \textit{hausdorff distance}:
\begin{def*}
	Let \((X, d)\) be a copmlete metric space. Then, the \textbf{Hausdorff Distance} between sets A and B in \(\mathcal{H}(X)\) is defined by
	\[
		h(A, B) = d(A, B)\vee d(B, A). \; \square
	\]
\end{def*}
\begin{prop*}
	The ``metric'' h defined in \(\mathcal{H}(X)\) as above is a metric.
\end{prop*}
\begin{proof*}
	For the first property, it follows from
	\[
		h(A, A) = d(A, A)\vee d(A, A) = 0 \vee 0 = 0;
	\]

	Symmetry of the metric is a result of
	\[
		\left.\begin{array}{ll}
			h(A, B) = d(A, B)\vee d(B,A) \\
			h(B, A) = d(B, A)\vee d(A, B)
		\end{array}\right\} \Rightarrow h(A, B) = h(B, A).
	\]

	The property that the metric is always non-negative has already been proven today, and the triangle inequality follows from reapplying the definition of many things already seen: pick \(a\in A\), such that for all c in C,
	\begin{align*}
		d(a, B) = \min\limits_{b\in B}(d(a, b)) & \leq \min\limits_{b\in B}(d(a, c) + d(c, b)) \\
		                                        & = d(a, c) + \min\limits_{b\in B}(c, b),
	\end{align*}
	so
	\[
		d(a, B)\leq \min\limits_{c\in C}(d(a, c)) + \max\limits_{b\in B}\{d(c, b):\; c\in C\}= d(a, C) + d(C, B).
	\]
	Because it holds for all a in A,
	\[
		d(A, B) \leq  d(A, C) + d(C, B).
	\]
	Similarly,
	\[
		d(B, A) \leq d(B, C) + d(C, A),
	\]
	form which we conclude that
	\begin{align*}
		h(A, B) = d(A, B)\vee d(B, A) & \leq d(B, C)\vee d(C, B) + d(A, C)\vee d(C, A) \\
		                              & = h(B, C) + h(A, C). \text{ \qedsymbol}
	\end{align*}
\end{proof*}
\begin{prop*}
	For every A, B, C and D in \(\mathcal{H}(X)\), we have
	\[
		h(A\cup B, C\cup D) \leq h(A, C) \vee h(B, D).
	\]
\end{prop*}
\begin{proof*}
	For the sets above, we know that
	\[
		h(A\cup B, C\cup D) = d(A\cup B, C\cup D)\vee d(C\cup D, A\cup B),
	\]
	thus, if the point x at which the maximum is attained is a point in A,
	\[
		d(A\cup B, C\cup D) = d(x, C\cup D) = \min(d(x, C), d(x, D)),
	\]
	so
	\[
		d(x, C\cup D) \leq d(x, C) \leq h(A, C) \leq h(A, C)\vee h(B, D).
	\]

	On the other hand, if the maximum is a point y in B, it follows similarly that
	\[
		d(y,C\cup D) \leq d(y, D) \leq h(B, D)\leq h(A, C)\vee h(B, D). \text{ \qedsymbol}
	\]
\end{proof*}

\begin{theorem*}
	Whenever the metric space \((X, d)\) is compact, \((\mathcal{H}(x), h)\) is also compact.
\end{theorem*}
\begin{proof*}
	As done a lot of the time, let's step down to points instead of sets by first considering a sequence of sets \(\{A_{k}\}_{k=1}^{\infty},\; A_{k}\in \mathcal{H}(X)\); then, pick \(x_{k}\in A_{k}\) to form a sequence of points \(\{x_{k}\}_{k=1}^{\infty}\), and by compactness of X, this sequence must have a subsequence with a limit point in X, say \(\{x_{k_{j}}\}_{j=1}^{\infty}\).

	We want to build a set A such that
	\[
		A = \{y\in X:\; \exists x_{k_{j}}\in A_{k_{j}},\; x_{k_{j}}\to y\}.
	\]

	First, we will prove that A is compact and nonempty; as a matter of fact, we know that it is not empty because, by compactness of X, every sequence has a converging subsequence in X, such that its limit will be a point in A by the way this set was taken.
	With respect to compactness, A must be closed, which means it must contain all of its limit points; thus, take a sequence \(\{a_{n}\}_{n=1}^{\infty}\) of elements of A, and assume it has a limit denoted by a.
	We are going to work with the limit of this sequence and a sequence we built from the subsequences of \(\{x_{k}\}\), then show they coincide.

	Pick \(\varepsilon  > 0 \), to which we can apply the definition of convergence to find a natural number m such that if \(n > m\), then
	\[
		d(a_{n}, y) \leq \frac{\varepsilon }{2}.
	\]
	Since \(a_{n}\) belong to A for all n's, there is a sequence \(\{x_{k}\}_{k=1}^{\infty}\), \(x_{k}\) in \(A_{k}\), such that for some \(k\geq n\),
	\[
		d(x_{k}, a_{n}) < \frac{\varepsilon }{2}.
	\]
	Hence,
	\[
		d(x_{k}, y) < d(x_{k}, a_{n}) + d(a_{n}, y) < \frac{\varepsilon }{2} + \frac{\varepsilon }{2},
	\]
	which means \(x_{k}\) converges to y, which, by definition, must belong to A.

	Finally, we will show that the initially chosen sequence \(\{A_{k}\}\) satisfies
	\[
		h(A_{k}, A) \to 0.
	\]
	Suppose, by means of contradiction, that it is not the case, so that
	\[
		\lim_{k\to \infty}h(A_{k}, A) > 0.
	\]
	By construction, there must be a subsequence \(A_{k_{j}}\) and \(\varepsilon  > 0\) such that, for all j,
	\[
		h(A_{k_{j}}, A)\geq \varepsilon;
	\]
	consequently, for any subsequence \(\{x_{k_{j}}\}\in A_{k_{j}}\), we must also have
	\[
		d(x_{k_{j}}, A)\geq \varepsilon,
	\]
	but because X is compact, \(\{x_{k_{j}}\}\) must have a subsequence with a limit point z in X, which means z must be in A so that \(d(z, A) = 0\). However, since \(d(x_{k_{j}}, A)\geq \varepsilon \) for all j, and in particular \(d(z, A)\geq \varepsilon \) -- we have reached a contradiction.

	Therefore,
	\[
		\lim_{k\to \infty}h(A_{k}, A) = 0. \text{ \qedsymbol}
	\]
\end{proof*}

\begin{tcolorbox}[
		skin=enhanced,
		title=Observation,
		fonttitle=\bfseries,
		colframe=black,
		colbacktitle=cyan!75!white,
		colback=cyan!15,
		colbacklower=black,
		coltitle=black,
		drop fuzzy shadow,
		%drop large lifted shadow
	]
	If we were to use the supremum and the infimum instead of the maximum and the minimum when defining \((\mathcal{H}(X), h)\) and its steps, then we would not get a metric in most cases.
	One of the problems is with the property \(h(A ,A) = 0\) and with \(h(A, B)\neq 0\) when \(A\neq B\); for instance, consider the sets \(A = (0, 1),\; B = [0,1]\), so that
	\[
		\tilde{d}(A, B) = \sup_{}\{\tilde{d}(x, B):\; x\in A\} = 0,
	\]
	because every point in the interior of A is also in B, but the distance between the boundaries is also 0, contradicting the definition of metric.

	Another counterexample is to consider \(\mathbb{R}\) and \(A = \{0\}\), in which case \linebreak \(h(A, R) = \infty\).
\end{tcolorbox}


\end{document}
